% Paper 2, §1 — Introduction
% Pass-C synthesis: previews §2-§14, all of which (except §2/§13/§14) are written.
% Anchors: kb/notes/training/, kb/notes/post-training/, brainstorm in
%          theory/series/SHAPE-C-decision.md §"Paper 2".

\section{Introduction}
\label{sec:introduction}

Training a frontier language model is not the optimization of a single
objective. It is a pipeline of objectives, each with its own loss form,
data distribution, optimizer, KL anchor, and stopping criterion. The
pre-training cross-entropy that fits a base model on $\sim\!15$T tokens
\citep[\S 3.2]{deepseek-v3} is not the same kind of object as the
classification loss on a curated $\sim\!10^4$-pair preference set
\citep{rafailov2023-dpo}, which is not the same kind of object as the
sparse $\{0,1\}$ verifier signal that drives a $30$K-sample GRPO run on
math contest problems
\citep{shao2024,deepseek-r1}. They live on different data, optimize
different functionals, anchor against different reference policies, and
stop on different criteria. Treating ``training'' as one optimization
collapses these distinctions and produces a class of empirics that has
genuinely confused the field for half a decade: the
Kaplan-vs-Chinchilla compute-allocation flip
\citep{kaplan2020,hoffmann2022-chinchilla}, the
DPO-vs-PPO-vs-GRPO debate, and the recurrent claim that ``RLVR creates
new reasoning capability'' are all, on inspection, artefacts of mixing
two stages' invariants. The thesis of this paper is that they
disentangle once the pipeline is named:
\begin{enumerate}
\item \textbf{Pre-training data.} What goes into the corpus, how it is
  filtered and mixed, how many tokens per parameter
  (\cref{sec:pre-training-data}).
\item \textbf{Scaling and optimization.} How loss falls with $(N, D, C)$
  and how to spend a fixed compute budget; what optimizer spends those
  FLOPs (\cref{sec:scaling-laws}, \cref{sec:optimization}).
\item \textbf{Distributed and mixed-precision scaffolding.} How the
  training step is laid out across thousands of accelerators, and at
  what numerical precision (\cref{sec:distributed-training},
  \cref{sec:mixed-precision}).
\item \textbf{Supervised fine-tuning.} The first post-training stage:
  cross-entropy on curated prompt-response pairs, producing the
  instruction-following starting point $\pi^{\mathrm{SFT}}$
  (\cref{sec:sft}).
\item \textbf{Preference-RL.} RLHF, RLAIF / Constitutional AI, DPO and
  the offline-preference variant family: aligning the policy against
  pairwise preference data, with a KL anchor to $\pi^{\mathrm{SFT}}$
  (\cref{sec:rlhf}, \cref{sec:dpo}, \cref{sec:rlaif-cai}).
\item \textbf{Reasoning-RL and adaptation.} RL on verifiable rewards
  (RLVR / GRPO / DAPO) for math, code, and proof, plus the weight-space
  operations (LoRA / QLoRA / model soups / TIES) that compose
  adapt-and-merge late in the pipeline
  (\cref{sec:rlvr-grpo}, \cref{sec:adaptation-and-merging}).
\end{enumerate}
Each stage has its own canonical lineage, its own load-bearing
variables, its own resolved questions, and its own live frontiers. The
pipeline composes; the stages are not interchangeable; and the
``confused empirics'' resolve as soon as we identify which stage's
invariant is being measured.

\subsection{What the pipeline framing buys}
\label{sec:intro-buy}

Three persistent debates in the 2020--2026 literature are stage-
boundary artefacts. We name them here because they motivate the rest of
the paper.

\paragraph{Kaplan vs. Chinchilla.} \citet{kaplan2020} fitted three
power laws on small-to-medium decoder-only LMs and recommended scaling
$N$ much faster than $D$ at fixed compute --- ``train very large models,
stop well short of convergence.'' \citet{hoffmann2022-chinchilla}
re-fit on a re-instrumented set of $400+$ runs whose LR schedules
matched their actual token horizons (Kaplan's cosine schedule decayed
to a fixed horizon biased short-data runs upward), found
$N \propto C^{0.5},\, D \propto C^{0.5}$, and produced the famous
$D \approx 20N$ rule along with the Chinchilla-70B-beats-Gopher-280B
headline at matched FLOPs. The two papers' \emph{functional forms}
agree; the \emph{compute-allocation prescription} flipped on the
strength of one methodological correction. The 2024--2026 frontier
sits an order of magnitude beyond Chinchilla on tokens-per-parameter
(Llama-3 8B at $D/N \approx 1875$, two orders past the
Chinchilla-optimal point), which is \emph{not} a refutation of
Chinchilla but a re-objective: pre-training cost has been displaced by
inference cost as the dominant economic axis
\citep[\S 3.2]{deepseek-v3}. \cref{sec:scaling-laws} treats this
in full.

\paragraph{RLHF vs. DPO vs. PPO.} \citet{rafailov2023-dpo} read the
RLHF closed-form RL optimum
$\pi^*(y \mid x) \propto \pi^{\mathrm{ref}}(y \mid x)\,\exp(r(x,y)/\beta)$
not as the \emph{target} of a PPO loop \citep{ouyang2022-instructgpt}
but as the \emph{definition} of a reward that can be inverted ---
yielding a closed-form classification loss on the same preference data,
without an explicit reward model, without on-policy sampling, without
a value head. Same fixed point, two-orders-cheaper per step. The 2023
narrative ``DPO replaces PPO'' was overstated: at frontier scale the
two coexist (Llama-3 and T\"ulu-3 ship DPO; DeepMind / Gemma-3 keep
PPO with weight-averaged reward models), and \citet{meng2024-simpo}
reframes the DPO family as a small organized lattice with three
knobs (loss shape, reference dependence, length normalization).
\cref{sec:rlhf} and \cref{sec:dpo} treat this in full.

\paragraph{``RLVR creates new capability.'' } The DeepSeek-R1
\citep{deepseek-r1} headline --- pass@1 on AIME 2024 climbing from
$15.6\%$ to $77.9\%$ over a single GRPO run, with the model emitting a
\texttt{wait}-token ``aha-moment'' partway through --- was widely read
as RL conjuring reasoning capability out of a base LM. The 2025
follow-up literature, including a NeurIPS 2025 oral, substantially
narrowed that claim: RLVR sharpens pass@1 on problems already solvable
at some pass@$K$, but does not expand the set of problems solvable at
any $K$; the capability ceiling is set by pre-training. RLVR is
sampling-efficient sharpening, not capability expansion --- a
distinction that only comes into focus once the stages are named.
\cref{sec:rlvr-grpo} treats this in full and \cref{sec:contradictions}
returns to the live disagreement.

These three are the load-bearing examples. Each is a stage-boundary
collision: a quantity measured at one stage, an invariant claimed at
another. The pipeline framing isolates them.

\begin{intuition}
A useful mental picture: each stage is a distinct
\emph{KL-constrained reweighting} of the model's output distribution.
Pre-training fits an unconstrained $\pi^{\mathrm{base}}$ to the corpus
distribution. SFT performs a forward-KL projection of $\pi^{\mathrm{base}}$
onto a demonstration distribution. RLHF / DPO / RLAIF perform a reverse-KL
reweighting of $\pi^{\mathrm{SFT}}$ toward higher reward, with the anchor
$\beta\,\KL(\pi_\theta \|\, \pi^{\mathrm{SFT}})$ controlling how far the
policy may drift. RLVR performs the same reverse-KL reweighting with the
reward replaced by a verifier indicator. Returning to canonical form: the
post-training stages all share the closed-form RL optimum
$\pi^*(y|x) \propto \pi^{\mathrm{ref}}(y|x)\,\exp(r(x,y)/\beta)$
(\cref{sec:rlhf}); only $r$ and $\pi^{\mathrm{ref}}$ change between
RLHF, RLAIF, DPO, and RLVR. The pipeline is a \emph{chain} of such
reweightings; mixing two stages' references and rewards is what
produces the confused empirics.
\end{intuition}

\subsection{Why six stages, why now}
\label{sec:intro-six}

The six-stage decomposition is not a logical necessity. It is the
empirical shape of frontier-2026 recipes once we read across the
disclosed stack: DeepSeek-V3 / R1, Llama 3 / 4, Qwen 2.5 / 3, Phi-4,
T\"ulu 3, OLMo 2, Mistral Large, Gemma 3. Every disclosed pipeline
since 2024 contains all six stages, in the same order, with the same
inputs and outputs at each stage boundary. This convergence is recent:
GPT-3 \citep{kaplan2020} shipped without stages 4--6; InstructGPT
\citep{ouyang2022-instructgpt} added stages 4--5 (SFT and PPO-RLHF);
the RLVR / GRPO / DAPO line \citep{shao2024,deepseek-r1,dapo2025} added
stage 6's reasoning-RL leg; the merging line
\citep{wortsman2022-model-soups,ilharco2022-task-vectors,yadav2023-ties-merging}
added stage 6's adaptation leg. The 2026 question --- \emph{is the
six-stage decomposition canonical, or recipe-zoo bias from the same few
labs reading each other's papers?} --- is open, and we treat it
explicitly in \cref{sec:frontier-snapshot} after the body of the paper
has put each stage's load-bearing variables on the page.

\subsection{Notation and scope}
\label{sec:intro-notation}

Paper 2 assumes Paper 1's residual-stream / FFN / attention notation.
The Transformer block, with its alternating attention and feed-forward
sub-layers around residual connections, is taken as given
\citep{vaswani2017}; the parameter count $N$, sequence length $S$,
batch size $B$, model dimension $D$, head count $H$, head dimension
$d_h$, vocabulary size $V$, and layer count $L$ from Paper 1's
preamble are reused without redefinition. New to this paper: $D$ also
denotes the training-token count when paired with $N$ (the scaling-law
context disambiguates), $C \approx 6ND$ is total non-embedding training
compute in FLOPs \citep[\S 2.1]{kaplan2020}, and the post-training
stages introduce $\pi_\theta$ for the trainable policy,
$\pi^{\mathrm{ref}}$ for whichever reference is anchoring it,
$r_\phi(x,y)$ or $R(x,y)$ for learned and verifiable rewards, and
$\beta$ for the KL coefficient.

We hand off a single forward thread to Paper 3:
\cref{sec:rlvr-grpo}'s GRPO and DAPO machinery is the \emph{training
leg} of Paper 3's compute / search / verification triangle for
reasoning. The \emph{inference} side --- chain-of-thought sampling,
self-consistency, process reward models, MCTS-style tree search at
generation time --- is treated there. Paper 4 (interpretability) and
Paper 5 (evaluation) reach back into this paper for, respectively, the
training-time activations that probes are built on and the alignment
surface that RLHF / CAI shapes.

What this paper deliberately does \emph{not} cover. (i) Architectural
choices --- attention variants, MoE, normalization placement,
positional encodings --- live in Paper 1, even though MoE intersects
training (\cref{sec:distributed-training} on EP, \cref{sec:mixed-precision}
on FP8 expert-collapse) and is cited from there. (ii) Inference-time
techniques --- speculative decoding, KV-cache compression, quantized
serving --- live in Paper 1's inference section. (iii) Evaluation
methodology --- contamination, benchmark validity, alignment-faking
detection --- lives in Paper 5; this paper takes \emph{loss}, \emph{pass@1},
and \emph{preference win-rate} as legitimate signals and defers
benchmark critique to Paper 5.

\subsection{Roadmap}
\label{sec:intro-roadmap}

The body of this paper walks the six stages in order, then closes with
a frontier snapshot and a contradictions chapter.

\begin{description}
\item[\cref{sec:pre-training-data} --- Pre-training data.] Curate
  $\to$ ablate $\to$ mix. The FineWeb pipeline end-to-end (96 CC
  snapshots, MinHash 14$\times$8, custom quantile-derived filters), the
  Data Mixing Laws functional form, and the FineWeb-Edu
  distill-then-classify pattern. Token-budget choice ($D \approx 20N$
  Chinchilla, $D/N \approx 1875$ Llama-3 over-train) sets up
  \cref{sec:scaling-laws}.
\item[\cref{sec:scaling-laws} --- Scaling laws.] Kaplan's three power
  laws and joint $L(N, D)$; Chinchilla's matched-LR-horizon re-fit
  yielding $N \propto C^{0.5},\, D \propto C^{0.5}$; the post-Chinchilla
  industry move to over-train for inference economics; \textmu-Transfer
  as the HP-transfer technology that makes any of these recipes tunable
  at frontier scale.
\item[\cref{sec:optimization} --- Optimization.] AdamW as the
  production default; Lion and Sophia as training-loss-only speedups;
  Muon's Newton--Schulz orthogonalization with a verified
  $\sim\!2\times$ compute-efficiency lift at $16$B-MoE / $5.7$T tokens;
  the schedule axis (cosine $\to$ Warmup--Stable--Decay $\to$
  Schedule-Free).
\item[\cref{sec:distributed-training} --- Distributed training.] Five
  composable axes on a single \texttt{DeviceMesh}:
  $G = d \cdot t \cdot c \cdot p \cdot e$ for DP $\times$ TP $\times$
  CP $\times$ PP $\times$ EP. Per-axis cost analysis, DeepSeek-V3's
  TP-free plan as case study, TorchTitan's 4D composability,
  Distributed Muon.
\item[\cref{sec:mixed-precision} --- Mixed precision and stability.]
  Three independent precision choices (storage / compute /
  communication) on a master-FP32 backbone; the format zoo
  (FP32 / FP16 / BF16 / FP8 / NVFP4 / BitNet); DeepSeek-V3's
  fine-grained FP8 recipe with zero irrecoverable loss spikes over
  $14.8$T tokens; loss-spike taxonomy and mitigations.
\item[\cref{sec:sft} --- Supervised fine-tuning.] Cross-entropy with
  prompt-masking; the data-not-loss thesis; lineage from FLAN through
  InstructGPT to T\"ulu-3, MAGPIE, R1-Distill, and s1; multi-turn
  loss-mask conventions; why SFT precedes RL.
\item[\cref{sec:rlhf} --- RLHF.] InstructGPT's three-stage pipeline;
  Bradley--Terry reward modelling; PPO with KL anchor and the
  closed-form RL optimum; the four-model PPO infrastructure;
  reward-model variants (ensembles, WARM/WARP).
\item[\cref{sec:dpo} --- DPO and offline preference.] The
  reparameterization that turns the RLHF optimum into a classification
  loss on the same preference data; the variant lattice (IPO, KTO,
  ORPO, SimPO, CPO) along three knobs.
\item[\cref{sec:rlaif-cai} --- RLAIF and Constitutional AI.] The same
  RLHF machinery with an AI labeler; CAI's two-stage pipeline (SL-CAI
  $\to$ RL-CAI); the constitution as the \emph{only} human-authored
  alignment input; non-evasion behavior and its dual-use surface.
\item[\cref{sec:rlvr-grpo} --- RLVR and GRPO.] Programmatic verifiers
  in place of learned reward models; group-mean baselines in place of
  co-sized value models; DeepSeek-R1's GRPO trajectory and the DAPO
  improvements (Clip-Higher, Dynamic Sampling, Token-Level Loss,
  Overlong Reward Shaping). Cross-paper anchor for Paper~3.
\item[\cref{sec:adaptation-and-merging} --- Adaptation and merging.]
  LoRA / QLoRA as low-rank parameter-efficient fine-tuning; Model
  Soups, Task Vectors, TIES-Merging as weight-space operations; the
  linear-mode-connectivity story as the speculative theoretical
  underpinning.
\item[\cref{sec:frontier-snapshot} --- Frontier-2026 snapshot.] The
  six-stage table across the disclosed frontier (DeepSeek-V3 / R1,
  Llama 3 / 4, Qwen 2.5 / 3, Phi-4, T\"ulu 3, OLMo 2, Mistral Large,
  Gemma 3). Where labs agree, where they diverge, where they are
  opaque. Returns to the open question of whether the six-stage
  decomposition is canonical or recipe-zoo bias.
\item[\cref{sec:contradictions} --- Live contradictions.] Concentrated
  treatment of the \texttt{[CONTRADICTION]} markers raised throughout
  the body: Kaplan-vs-Chinchilla, optimizer speedups, TP-yes-vs-TP-no,
  FP8 / NVFP4 / BitNet production status, SFT data quantity-vs-quality,
  reward-hacking, DPO-vs-PPO, CAI internalization, RLVR capability,
  merge guarantees, precision $\times$ token-count, \textmu-P at
  frontier scale.
\end{description}

The closing thesis we will defend in \cref{sec:contradictions}: the
pipeline's open questions are \emph{localized}. Most live at stage
boundaries (precision $\times$ optimizer; token-count $\times$
quantization; preference $\times$ verifier) rather than at the
pipeline's structure. The structure --- pre-train $\to$ scale $\to$
optimize $\to$ distribute $\to$ precision $\to$ SFT $\to$ preference
$\to$ reasoning $\to$ merge --- is now stable enough to write down.
The empirical envelope around each stage's variables is what
2027-class results will close.
